\documentclass[12pt,a4paper]{article}

\usepackage[utf8]{inputenc}
\usepackage[T1]{fontenc}
\usepackage[brazilian]{babel}
\usepackage{geometry}
\usepackage{setspace}
\usepackage{titlesec}
\usepackage{hyperref}
\usepackage{xcolor}
\usepackage{tcolorbox}
\usepackage{fancyhdr}
\usepackage{booktabs}
\usepackage{tikz}
\usepackage{enumitem}
\usepackage{parskip}
\usepackage{amsmath}
\usepackage{graphicx}
\usepackage{float}
\usepackage{caption}
\usetikzlibrary{positioning}

\geometry{left=2.5cm, right=2.5cm, top=2.5cm, bottom=2.5cm}
\setstretch{1.10}
\emergencystretch=2em

% -- Cores ---------------------------------------------------------------
\definecolor{azul}{HTML}{1A3A6B}
\definecolor{azulclaro}{HTML}{D5E8F4}
\definecolor{cinzaborda}{HTML}{CCCCCC}

% -- Titulos -------------------------------------------------------------
\titleformat{\section}
  {\bfseries\large\color{azul}}{\thesection.}{0.5em}{}[\vspace{-2pt}\rule{\linewidth}{0.8pt}]
\titleformat{\subsection}{\bfseries\normalsize}{\thesubsection.}{0.5em}{}
\titlespacing{\section}{0pt}{12pt}{4pt}
\titlespacing{\subsection}{0pt}{8pt}{2pt}

% -- Caixas ---------------------------------------------------------------
\tcbuselibrary{skins, breakable}
\newtcolorbox{resultbox}{
  enhanced, colback=azul!7!white, colframe=azul!60!white,
  boxrule=0.6pt, arc=4pt, left=8pt, right=8pt, top=4pt, bottom=4pt,
  breakable,
}

% -- Cabecalho/rodape ----------------------------------------------------
\pagestyle{fancy}
\fancyhf{}
\renewcommand{\headrulewidth}{0pt}
\renewcommand{\footrulewidth}{0.4pt}
\fancyfoot[L]{\small\color{gray} SI405 -- Adapter e Decorator}
\fancyfoot[R]{\small\color{gray} 2026}
\fancyfoot[C]{\small\color{gray}\thepage}

\hypersetup{colorlinks=true, linkcolor=azul, urlcolor=azul}
\captionsetup{font=small, labelfont=bf}
\graphicspath{{si405_report_assets/}{./}}

\begin{document}
\thispagestyle{empty}

% ========================================================================
% CAPA
% ========================================================================
\begin{titlepage}
  \centering

  {\Large\bfseries Universidade Estadual de Campinas}\\
  {\large Faculdade de Tecnologia -- FT}\\[1.5cm]

  \IfFileExists{af-logo-ft-reduzido-pos.png}{
    \includegraphics[width=0.25\textwidth]{af-logo-ft-reduzido-pos.png}\\[0.5cm]
  }{
    \begin{tikzpicture}
      \draw[azul, line width=1.2pt, rounded corners=6pt] (0,0) rectangle (4.2,1.5);
      \node[azul] at (2.1,0.75) {\Large\bfseries FT -- UNICAMP};
    \end{tikzpicture}\\[0.5cm]
  }

  {\LARGE\bfseries\color{azul} Relatório Técnico}\\
  {\Large SI405 -- Padrões Estruturais: Adapter e Decorator}\\[0.2cm]

  \vspace{0.5cm}
  \begin{tikzpicture}
    \draw[azul, line width=1.8pt] (0,0) -- (0.96\textwidth,0);
  \end{tikzpicture}
  \vspace{0.5cm}

  \vfill

  \begin{flushright}
    \textbf{Discentes:}\\[0.3em]
    Daniel Aniceto Rosell -- RA 283988\\[0.6em]
    Davie Schimidt Fonseca -- RA 259908\\[0.6em]
    Hugo Strassa -- RA 246710\\[0.6em]
    Gabriel Sorensen M. Traina -- RA 283997\\[0.6em]
    Kaue Samuel Oliveira da Silva -- RA 178449\\[0.6em]
    Kauã Henrique da Silva Andrade -- RA 246165\\[0.8em]
  \end{flushright}

  \vspace{1cm}
  \textbf{Limeira -- SP}\\
  {\large 2026}
\end{titlepage}

\setcounter{page}{1}

% ========================================================================

\section{Introdução}

Este relatório apresenta a solução desenvolvida para a atividade prática de SI405 sobre padrões estruturais, com foco nos padrões \textit{Adapter} e \textit{Decorator}. O sistema-base representa uma aplicação corporativa de cobrança de pedidos, originalmente capaz de processar pagamentos por boleto, Pix e cartão de crédito, além de aplicar ajustes de valor como desconto, juros, taxa internacional e seguro.

A implementação inicial estava funcional, porém apresentava problemas arquiteturais importantes. A integração com gateways de pagamento era feita diretamente no serviço principal, inclusive com código específico para bibliotecas externas. Além disso, os ajustes de valor eram controlados por parâmetros booleanos, o que dificultava a evolução da aplicação quando novas taxas precisassem ser adicionadas.

\begin{resultbox}
\textbf{Repositório GitHub da solução:} \url{https://github.com/SorensenG/si405-padroes-estruturais-grupo}
\end{resultbox}

A solução preservou o projeto original, manteve os testes existentes funcionando e adicionou suporte à forma de pagamento Carteira Digital e aos novos ajustes exigidos no enunciado.

\section{Análise dos principais problemas arquiteturais}

O primeiro problema identificado foi o acoplamento do serviço de cobrança às implementações concretas dos gateways. A classe \textit{CobrancaService} conhecia diretamente o gateway interno, o SDK externo \textit{PaySecureGateway}, o formato dos dados exigidos por esse SDK e o tratamento de exceções específico da integração. Dessa forma, uma classe de regra de aplicação passou a depender de detalhes de infraestrutura.

Outro problema estava na duplicação da lógica de seleção e conversão de gateways. Os métodos \textit{cobrar} e \textit{cobrarEmLote} repetiam a mesma estrutura condicional para boleto, Pix e cartão de crédito. Assim, qualquer nova forma de pagamento exigiria alterações em mais de um ponto, aumentando o risco de inconsistência.

No cálculo de valor final, a aplicação usava parâmetros booleanos para indicar quais ajustes deveriam ser aplicados. Essa solução funcionava para poucos casos, mas não era extensível. Cada novo ajuste exigiria alterar a assinatura de métodos, adicionar novos condicionais e espalhar a regra pelo serviço.

\begin{resultbox}
\textbf{Problema central:} o sistema funcionava, mas misturava seleção de gateway, adaptação de SDK externo, cálculo de valor e composição de ajustes dentro do mesmo fluxo de serviço.
\end{resultbox}

\section{Justificativa das decisões arquiteturais adotadas}

A solução adotada foi uma refatoração localizada, mantendo o código-base fornecido e evitando substituir o sistema por uma implementação nova. A prioridade foi preservar o comportamento existente, reduzir acoplamento e criar pontos claros de extensão.

Para os meios de pagamento, foi utilizado o padrão \textit{Adapter}. O serviço de cobrança passou a depender de uma interface comum de gateway, enquanto cada integração concreta ficou encapsulada em uma classe adaptadora. Essa decisão atende diretamente à restrição de não modificar classes do pacote externo, como \textit{PaySecureGateway} e \textit{WalletPaySDK}.

Para os ajustes de valor, foi utilizado o padrão \textit{Decorator}. Cada ajuste passou a ser representado por um objeto independente, capaz de envolver o valor calculado anteriormente e acrescentar sua regra. Com isso, a aplicação pode combinar ajustes em qualquer quantidade e ordem sem aumentar a assinatura principal de cobrança.

\section{Aplicação do padrão Adapter}

O padrão \textit{Adapter} foi aplicado por meio da interface \textit{GatewayCobranca}. Essa interface define uma operação comum de cobrança, recebendo o pedido, o valor final e a forma de pagamento. O \textit{CobrancaService} deixa de conhecer diretamente os SDKs e passa a depender apenas dessa abstração.

Foram criados três adaptadores principais. O \textit{GatewayPagamentoInternoAdapter} adapta o gateway interno usado por boleto e Pix. O \textit{PaySecureGatewayAdapter} adapta o SDK externo de cartão de crédito. O \textit{WalletPayGatewayAdapter} integra o novo SDK de Carteira Digital solicitado no enunciado.

No caso do \textit{PaySecureGateway}, o adapter centraliza a criação do mapa de dados esperado pelo SDK, incluindo identificador do pedido, nome do cliente, valor e moeda. Também converte o código de status externo para os status internos \textit{APROVADA} e \textit{RECUSADA}, além de tratar a exceção de indisponibilidade do gateway.

No caso do \textit{WalletPaySDK}, o adapter converte o valor de reais para centavos, cria a requisição exigida pelo SDK e traduz os status externos. O status \textit{CONFIRMED} é tratado como aprovado; os status \textit{DECLINED} e \textit{FAILED} são tratados como recusados.

\begin{table}[H]
\centering
\small
\begin{tabular}{p{0.30\textwidth}p{0.52\textwidth}}
\toprule
\textbf{Componente} & \textbf{Responsabilidade} \\
\midrule
\textit{GatewayCobranca} & Interface comum usada pelo serviço de cobrança. \\
\textit{GatewayPagamento}\newline\textit{InternoAdapter} & Adapta boleto e Pix para a interface comum. \\
\textit{PaySecureGatewayAdapter} & Isola a integração com o SDK externo de cartão. \\
\textit{WalletPayGatewayAdapter} & Integra a nova forma de pagamento Carteira Digital. \\
\textit{GatewayCobrancaResolver} & Centraliza a escolha do adapter adequado para cada forma de pagamento. \\
\bottomrule
\end{tabular}
\caption{Organização da solução baseada em Adapter.}
\end{table}

\section{Aplicação do padrão Decorator}

O padrão \textit{Decorator} foi aplicado no cálculo do valor final da cobrança. O componente base é representado por \textit{ValorCobranca}, e a classe \textit{ValorBaseCobranca} encapsula o valor inicial do pedido. A classe abstrata \textit{AjusteValorDecorator} serve como base para os decorators concretos.

Os ajustes já existentes foram preservados nas classes \textit{DescontoFidelidadeDecorator}, \textit{JurosParcelamentoDecorator}, \textit{TaxaOperacaoInternacionalDecorator} e \textit{SeguroTransacaoDecorator}. As regras originais foram mantidas: desconto de 5\%, juros de 2,99\%, taxa internacional de 5\% e seguro fixo de R\$ 4,90.

Também foram adicionados os novos ajustes solicitados: \textit{TaxaAntecipacaoRecebiveisDecorator}, com taxa de 1,5\%, e \textit{TaxaEmissaoNotaFiscalDecorator}, com valor fixo de R\$ 2,50. O serviço passou a aceitar uma lista de ajustes, permitindo montar a cadeia de decorators dinamicamente.

\begin{resultbox}
\textbf{Compatibilidade:} os métodos antigos com parâmetros booleanos foram mantidos como camada de compatibilidade, mas internamente delegam para a nova composição orientada a objetos.
\end{resultbox}

\begin{table}[H]
\centering
\small
\begin{tabular}{p{0.38\textwidth}p{0.44\textwidth}}
\toprule
\textbf{Decorator} & \textbf{Regra aplicada} \\
\midrule
Desconto de fidelidade & Reduz 5\% do valor. \\
Juros de parcelamento & Acrescenta 2,99\% ao valor. \\
Taxa internacional & Acrescenta 5\% ao valor. \\
Seguro de transação & Acrescenta R\$ 4,90. \\
Antecipação de recebíveis & Acrescenta 1,5\% ao valor. \\
Emissão de nota fiscal & Acrescenta R\$ 2,50. \\
\bottomrule
\end{tabular}
\caption{Ajustes representados como decorators.}
\end{table}

\section{Impactos da refatoração}

A refatoração reduziu o acoplamento entre o serviço de cobrança e os gateways concretos. Para adicionar uma nova forma de pagamento, basta criar um novo adapter e registrá-lo no resolvedor. O serviço principal permanece estável e não precisa conhecer os detalhes de cada SDK.

A duplicação entre cobrança individual e cobrança em lote também foi reduzida. Os dois fluxos reutilizam a mesma lógica de cálculo e a mesma resolução de gateway, evitando que mudanças futuras precisem ser repetidas manualmente em mais de um método.

Com o uso de \textit{Decorator}, os ajustes de valor ficaram independentes, combináveis e testáveis isoladamente. A criação de novas taxas não exige novos parâmetros booleanos nem condicionais no cálculo central. Essa organização favorece manutenção e extensibilidade, além de demonstrar o princípio estudado em aula: preferir composição sobre estruturas rígidas de decisão.

\section{Testes executados}

Os testes existentes foram mantidos e novos cenários foram adicionados para cobrir os requisitos do enunciado. Foram criados testes para o adapter do \textit{PaySecureGateway}, para o adapter do \textit{WalletPaySDK}, para os novos decorators e para combinações de múltiplos decorators em ordens diferentes.

\begin{resultbox}
\textbf{Resultado da validação:} a execução de \texttt{mvn test} finalizou com 34 testes executados, 0 falhas e 0 erros. Também foi executado \texttt{mvn clean package -q} com sucesso.
\end{resultbox}

\begin{table}[H]
\centering
\small
\begin{tabular}{p{0.30\textwidth}p{0.52\textwidth}}
\toprule
\textbf{Categoria de teste} & \textbf{Cenários validados} \\
\midrule
Testes existentes & Pedido, gateway interno, PaySecure, WalletPay e serviço de cobrança. \\
Adapters & Aprovação e recusa em cartão e Carteira Digital. \\
Decorators novos & Antecipação de recebíveis e emissão de nota fiscal. \\
Composição & Combinação de múltiplos decorators em ordens distintas. \\
\bottomrule
\end{tabular}
\caption{Resumo dos testes automatizados.}
\end{table}

\end{document}
